\section{Wprowadzenie i motywacja}

\subsection{Kontekst problemu}
Aplikacje webowe są stale narażone na automatyczny ruch generowany przez boty:
masowe skanowanie zasobów, testowanie endpointów, tworzenie sztucznego ruchu
lub działania fraudowe. Klasyczne podejścia ochronne oparte wyłącznie o reguły
sieciowe i statyczne sygnatury przestają być wystarczające, ponieważ nowoczesne
boty potrafią imitować część zachowań użytkownika.

\subsection{Dlaczego behawioralna biometria myszy}
Trajektoria kursora, rytm kliknięć, dynamika scrollowania i przejścia focus/blur
niosą bogatą informację o sposobie interakcji człowieka z interfejsem.
W odróżnieniu od pojedynczego eventu, analiza krótkiego okna zachowania:

\begin{itemize}
  \item lepiej oddaje \textit{styl} interakcji,
  \item utrudnia prostym botom przejście niezauważonym,
  \item daje możliwość szybkiej reakcji ochronnej jeszcze w trakcie sesji.
\end{itemize}

\subsection{Cel projektu}
Celem MVP było zbudowanie kompletnego pipeline'u \textit{end-to-end}:
\begin{itemize}
  \item frontend zbiera telemetrykę zachowania i fingerprint środowiska,
  \item backend zapisuje sesje i surowe eventy,
  \item moduł cech buduje reprezentację okien zachowania,
  \item model klasyfikuje \texttt{human} vs \texttt{bot},
  \item silnik polityk zwraca akcję: \texttt{allow/observe/throttle/challenge/block}.
\end{itemize}

\subsection{Zakres i świadome uproszczenia MVP}
Wersja demonstracyjna została zaprojektowana profesjonalnie, ale celowo uproszczona:
\begin{itemize}
  \item inferencja modelu odbywa się w tym samym backendzie FastAPI (bez osobnego model service),
  \item predykcja działa na poziomie sesji i krótkich okien eventów,
  \item brak integracji z zewnętrznym WAF i blokadą IP (reakcja sesyjna),
  \item nacisk na czytelność kodu, odtwarzalność wyników i łatwą prezentację lokalną.
\end{itemize}
